\begin{abstract}
\noindent
A machine that makes many different parts needs many different tools, but it can
hold only a few of them at a time. Running a set of jobs therefore means stopping
to change tools, and each stop costs production time. The Job Sequencing and Tool
Switching Problem (SSP) asks for the order of jobs, and the tool loading at each
step, that makes the total number of tool changes as small as possible. Choosing
the loading for a fixed order is easy. Choosing the order is NP-hard, and it is
where all of the difficulty sits.

In practice the SSP is not attacked directly. Jobs are first grouped into as few
magazine-fitting batches as possible, and the batches are then sequenced. This
two-phase method is fast and is what industry uses, but no worst-case analysis of
it has been published. This report gives one. It first shows that the method
admits two different readings, which the literature has not separated: the cost of
a walk through one fixed magazine per group, and the cost the method actually
returns once its final step re-optimises the loading. The two differ, and they
differ on the instances most often used as examples. Separating them is what makes
the rest of the analysis correct. Within that setting, the optimum is shown to
equal the cheapest walk over \emph{all} groupings, so the two-phase gap is exactly
the price of insisting on the fewest groups; classes of instances with zero gap are
identified for every capacity; upper bounds on the gap are proved for every
instance; and the extremal behaviour is worked out at three groups. Positive gaps
turn out to be rare, small, and confined to instances whose magazine is large
relative to the jobs, which agrees with what the one published industrial study
reports. A parameterised cost model shows the SSP and the grouping problem to be
the two ends of a single spectrum, with the transition governed by the gap itself.

On the exact side the obstacle is a single quantity. The strongest compact
formulation in the literature has a linear relaxation equal to the number of tools
minus the magazine capacity, which is an elementary counting bound. Two methods are
built to get past it. A branch-and-Benders-cut algorithm places the ordering in a
master problem and settles the loading exactly with a polynomial oracle, so it never
forms a weak relaxation. Position-indexed formulations replace the configuration
walk by a fixed row set of polynomial size, and one of them, the
position--transition formulation, has a linear relaxation that provably exceeds the
counting bound.

Both were implemented and benchmarked against three reimplemented compact models on
$1{,}301$ standard instances under a single one-hour limit. The campaign is still
running; what follows is the snapshot available at the time of writing, and every
comparison is made on instances that the configurations being compared both attempted.
The results are consistent, and negative in an informative way. The branch-and-Benders-cut solver
closes $94\%$ of the instances on which the counting bound is tight and $49\%$ of
those on which it is not. Where it fails, the incumbent is usually already optimal
and only the proof is missing. Fractional Benders cuts, which the earlier campaign
could not test because of an implementation defect, were generated in very large
numbers once the defect was corrected; they raised the dual bound on three
instances out of $1{,}238$ and cost sixty-nine solved instances. Of four
strengthenings, only the one that improves the incumbent helps. The evidence is
that a configuration-based exact method for the SSP is limited by the strength of
its bound and not by its implementation.
\end{abstract}
